\section{Results (Partial)}

\begin{table}[!t]
\caption{Evidence scope. Caps are assisted speeds in m/s; a cell is one condition and setting.}
\label{tab:evidence}
\centering\scriptsize\setlength{\tabcolsep}{3pt}
\begin{tabular}{@{}p{0.25\columnwidth}p{0.36\columnwidth}p{0.31\columnwidth}@{}}
\toprule
Evidence & Configuration & Sample / purpose \\
\midrule
HTTP integration & Uncoupled; cap 0.15 & 9 trials; reliability \\
Intent sweep & Uncoupled; cap 0.15; vary $h,\sigma$ & 400 reaches/setting \\
Coupling test & Off/on; $k=0.75$, cap 0.15 & 400 reaches per cell \\
Speed comparison & $k=0.75$, cap 0.15 vs.\ $k=1.5$, cap 0.6 & 120 reaches per cell \\
Headset sessions & $k=1.5$, cap 0.6 & 1 tester; 15 + 18 completed trials \\
\bottomrule
\end{tabular}
\end{table}

\subsection{Concurrent Control in Integration Runs}
The service and player completed nine scripted HTTP trials without refused or
lost signals, causal-motion violations, or data loss. At 72 frames/s,
observation-to-application latency had medians of 16--40\,ms and a 95th
percentile of 42\,ms.

\begin{figure}[!t]
\centering
% Completed session 2 only: six trials/condition, one developer.
% Read rows 3--5 (manual, nearest, intent) from the verified session extract.
\pgfplotstableread[col sep=comma]{figures/headset-sessions.csv}\headsetdata
\newcommand{\readheadset}[3]{%
  \pgfplotstablegetelem{#1}{#2}\of\headsetdata\let#3\pgfplotsretval}
\readheadset{3}{time_median_s}{\manualTime}
\readheadset{4}{time_median_s}{\nearestTime}
\readheadset{5}{time_median_s}{\intentTime}
\readheadset{4}{false_assist_episodes_total}{\nearestEpisodes}
\readheadset{5}{false_assist_episodes_total}{\intentEpisodes}
\readheadset{3}{depth_effort_ratings}{\manualEffort}
\readheadset{4}{depth_effort_ratings}{\nearestEffort}
\readheadset{5}{depth_effort_ratings}{\intentEffort}
\def\splitratings#1;#2\relax{\def\firstRating{#1}\def\secondRating{#2}}
{\footnotesize\textbf{Session 2: one developer}\par}\vspace{3pt}
\begin{tikzpicture}
\pgfplotsset{headsetpanel/.style={
  scale only axis, height=2.30cm, font=\footnotesize,
  ymin=-0.90, ymax=2.55, ytick={0,1,2},
  axis x line*=bottom, axis y line*=left,
  tick align=outside, xmajorgrids, grid style={gray!20},
  title style={align=center, font=\footnotesize},
  clip=false,
}, timepoint/.style={only marks, mark=|, mark size=4pt, ultra thick}}
\begin{axis}[headsetpanel, name=timeplot, width=2.10cm,
  xmin=0, xmax=21, xtick={0,10,20},
  yticklabels={Intent,Nearest,Manual},
  title={\textbf{Trial time}\\Seconds}]
% The trial CSV is ordered manual (0--5), nearest (6--11), intent (12--17).
\addplot[only marks, mark=o, mark size=1.4pt, gray!80!black,
  restrict expr to domain={\coordindex}{0:5}]
  table[col sep=comma, x=time_s, y=plot_y] {figures/headset-session2-trials.csv};
\addplot[only marks, mark=o, mark size=1.4pt, blue!65!black,
  restrict expr to domain={\coordindex}{6:11}]
  table[col sep=comma, x=time_s, y=plot_y] {figures/headset-session2-trials.csv};
\addplot[only marks, mark=o, mark size=1.4pt, orange!85!black,
  restrict expr to domain={\coordindex}{12:17}]
  table[col sep=comma, x=time_s, y=plot_y] {figures/headset-session2-trials.csv};
\addplot[timepoint, gray!80!black] coordinates
  {(\manualTime,2)};
\addplot[timepoint, blue!65!black] coordinates
  {(\nearestTime,1)};
\addplot[timepoint, orange!85!black] coordinates
  {(\intentTime,0)};
\node[font=\scriptsize\bfseries, anchor=north, inner sep=0.5pt]
  at (axis cs:\manualTime,1.65) {\pgfmathprintnumber[fixed,precision=1]{\manualTime}};
\node[font=\scriptsize\bfseries, anchor=north, inner sep=0.5pt]
  at (axis cs:\nearestTime,0.65) {\pgfmathprintnumber[fixed,precision=1]{\nearestTime}};
\node[font=\scriptsize\bfseries, anchor=north, inner sep=0.5pt]
  at (axis cs:\intentTime,-0.35) {\pgfmathprintnumber[fixed,precision=1]{\intentTime}};
\end{axis}
\begin{axis}[headsetpanel, name=effortplot, width=1.55cm,
  at={(timeplot.south east)}, anchor=south west, xshift=0.55cm,
  xmin=0.5, xmax=7.5, xtick={1,4,7}, yticklabels={,,},
  title={\textbf{Depth effort}\\Rating (1--7)}]
\expandafter\splitratings\manualEffort\relax
\addplot[only marks, mark=*, mark size=2pt, gray!80!black]
  coordinates {(\firstRating,2.10) (\secondRating,1.90)};
\expandafter\splitratings\nearestEffort\relax
\addplot[only marks, mark=*, mark size=2pt, blue!65!black]
  coordinates {(\firstRating,1.10) (\secondRating,0.90)};
\expandafter\splitratings\intentEffort\relax
\addplot[only marks, mark=*, mark size=2pt, orange!85!black]
  coordinates {(\firstRating,0.10) (\secondRating,-0.10)};
\node[font=\scriptsize, align=center, anchor=north, yshift=-15pt]
  at (rel axis cs:0.5,0) {1 = none\\7 = a lot};
\end{axis}
\begin{axis}[headsetpanel, width=1.80cm,
  at={(effortplot.south east)}, anchor=south west, xshift=0.55cm,
  xmin=0, xmax=12, xtick={0,5,10}, yticklabels={,,},
  title={\textbf{Wrong-target}\\State episodes},
  xbar, bar width=7pt, bar shift=0pt,
  nodes near coords, point meta=x,
  every node near coord/.append style={font=\scriptsize, text=black}]
\addplot[fill=blue!30, draw=blue!65!black] coordinates {(\nearestEpisodes,1)};
\addplot[fill=orange!40, draw=orange!85!black] coordinates {(\intentEpisodes,0)};
\node[font=\scriptsize] at (axis cs:5,2) {N/A};
\node[font=\scriptsize, align=center, anchor=north, yshift=-15pt]
  at (rel axis cs:0.5,0) {Engagement coverage\\Nearest 100\%; Intent 75\%};
\end{axis}
\end{tikzpicture}

\caption{Complete session 2; six trials/condition. Time: trials (circles),
medians (bars and labels). Effort: two block ratings. Marks are offset vertically.
Episodes are state-based totals; manual has no AI. Engagement is median trial
coverage. Lower counts accompany less engagement.}
\label{fig:headset}
\end{figure}

\subsection{Intent Inference and Assistance Availability}
Compared with nearest, the chosen setting ($h = 0.5$, $\sigma = 0.08$) about
halved false-assist time toward distractor P (0.511 vs.\ 1.006\,s per reach crossing P).
The cost was lower coverage (0.327 vs.\ 0.387) and a
longer median reach-start-to-first-correct-assistance time (2.01 vs.\ 1.67\,s;
among reaches receiving correct assistance). These parameters were selected
on the same sweep; the comparison is exploratory.

\subsection{Coordinating Human Motion and Assisted Depth}
During 1\,s hovers over P (75 per condition), depth travel \emph{after} the
0.15\,s grace was 110.5\,mm (nearest) and 123.7\,mm (intent) without coupling,
and zero with it. Coupled nearest still moved 21.7\,mm during the grace.
In 400 reaches/condition, coupling required extra scripted movement on
32\% (nearest) and 77\% (intent), versus 0\% and 9\% without coupling.

In the separate speed comparison, raising $k$ and the cap reduced extra-motion
rates from 30\% to 17\% (nearest) and 72\% to 53\% (intent).
Mean wrong-hold travel rose from 49.7 to 83.1\,mm per reach (nearest) and
18.2 to 33.1\,mm (intent): faster recovery increased misdirected travel.

\subsection{Headset Testing (Developer)}
After a pilot stopped at discomfort 4, interrupted session 1 completed
15 measured trials (M: 3, N: 6, I: 6), excluding three incomplete attempts.
Session 2 completed 18/18 (six/condition) plus three practice trials, without
data loss, causal-motion violations, or no-signal pauses. All measured
completions used the final configuration.

Median trial times (M/N/I) were 36.5/11.7/13.4\,s in session 1 and
11.2/9.8/11.4\,s in session 2 (Fig.~\ref{fig:headset}). Wrong-target assistance episodes (N/I,
session totals) were 21/7 and 9/2, respectively. In session 2, median
wrong-hold depth travel was 39.6/0\,mm per trial, but maximum single-reach
travel was 185.8/93.3\,mm. Extra motion was required on 0/24 nearest and
1/24 intent cued reaches.

The developer reported lower effort with assistance (1--2 versus 7),
discomfort 1, no clear nearest/intent difference, and no stereo artifacts
after fixes.
